\section{Appendix}

\subsection{Bench-test results template}

The following table is provided blank; it is filled in as each
bench test from Chapter~\ref{sec:hil} is executed. The
date and the operator initials are recorded together with each
pass/fail mark so that the same table also serves as the
verification log for the report.

\begin{table}[ht]
\centering
\caption{Bench-test results log (to be filled in).}
\small
\begin{tabular}{@{}lp{0.4\linewidth}llll@{}}
\toprule
\textbf{ID} & \textbf{Test} & \textbf{Date} & \textbf{Operator} & \textbf{Result} & \textbf{Notes} \\
\midrule
A & PWM waveform sign-off               &  &  &  &  \\
B & ADC calibration vs known voltage    &  &  &  &  \\
C & eFuse switching on real board       &  &  &  &  \\
D & State machine on real sensor source &  &  &  &  \\
E & MPPT against a fake solar source    &  &  &  &  \\
F & Buck transfer curve                 &  &  &  &  \\
G & INA226 sensor chip verification     &  &  &  &  \\
H & CHIPS communication stress          &  &  &  &  \\
I & Watchdog timer                      &  &  &  &  \\
\bottomrule
\end{tabular}
\end{table}

\subsection{Build commands and repository layout}

The firmware is built with the GNU ARM toolchain
(\texttt{arm-none-eabi-gcc} 12.2.1) and flashed with OpenOCD
or Microchip Studio. The two relevant build targets are:

\begin{lstlisting}[language=bash]
make BOARD=devboard-pds   flash   # Curiosity Nano dev board
make BOARD=mainboard-pds  flash   # Real PCU testing board V4.1
\end{lstlisting}

The repository layout is summarised in Section~\ref{sec:architecture}.2.
Every file inside \texttt{src-pds/} is described by its own
file-naming convention; every file inside \texttt{src/drivers/}
has a header comment indicating its build target (\texttt{devboard
only}, \texttt{mainboard only}, or \texttt{shared}).

\subsection{External references}

The most useful external references for the firmware are listed
below.

\begin{itemize}
\item Microchip SAMD21 datasheet (full chip and peripheral
reference).
\item Microchip ATSAMD21 Device Family Pack (DFP), v3.6.144
(register definitions and startup code used by the firmware).
\item EPC2152 datasheet (the GaN half-bridge driving the buck).
\item Texas Instruments INA226 datasheet (the digital
current/voltage monitors on the I\textsuperscript{2}C bus).
\item Texas Instruments TPS25940 datasheet (the eFuses on the
panel and battery rails).
\item NASA Power of 10 \cite{holzmann_powerof10_2006} and JPL
Institutional Coding Standard for C \cite{jpl_d60411_2009} ---
the rule sources for the project's coding conventions.
\end{itemize}
